\scanpage{21}
% L1–3
qui s'explicitent ainsi : \deleted{\unsure{soit}} $a\colon\widehat C\longrightarrow\widetilde C$ le foncteur ``faisceau associé'', donc $\varepsilon_C=a\,i_C$ ($i_C\colon C\longrightarrow\widehat C$) donc pour $u\in\underline{\mathrm{Hom}}(\unsure{\overline C|\widetilde C},E)$, $u\circ\varepsilon_C=(ua)i_C$, donc (\deleted{\unsure{compte tenu}} que $a$, donc $ua$, commute aux petites $\varinjlim$)

% L4
\[
\alpha'(\unsure{u})=u\circ a
\]

% L5
\deleted{Donc} D'autre part, on voit que pour $v\in\underline{\mathrm{Hom}}'(\widehat C,E)$ (correspondant à $v_0\colon C\longrightarrow E$)

% L6
\[
\beta'(v)=v\circ j_{\unsure{C|a}}=v\mid\widetilde C\qquad(\text{où }j\colon\widetilde C\longrightarrow\widehat C\text{ est l'inclusion})
\]

% L7
\[
\text{On a bien}\qquad\beta'\alpha'(u)=u\circ a\mid\widetilde C\simeq u
\]

% L8–12
Le fait que $\alpha'$ est pl. fidèle provient du fait que $\widetilde C$ est, moyennant $a$, une catégorie de fractions de $\widehat C$, de sorte que $\underline{\mathrm{Hom}}(\widetilde C,E)$ (sans $'$) \deleted{est} s'identifie à la sous catégorie pleine de $\underline{\mathrm{Hom}}(\widehat C,E)$ formée des foncteurs \inserted{$F$} qui sont tels que \deleted{\unsure{$F(u)$}} $u\in\mathrm{Fl}\,\widehat C$, $a(u)$ inversible $\Longrightarrow F(u)$ inversible. Est-il vrai
% L13–15
que dans cette identification, les foncteurs qui commutent aux $\varinjlim$ petites se correspondent ? Il reste à voir que si $\unsure{H|U|\mathcal U}\colon\widetilde C\longrightarrow E$ est tel que $H\circ a$ commute aux $\varinjlim$, alors
% L16–19
$H$ aussi, or on aura $H=(H\circ a)i$, et si on a des objets $X_\alpha$ de $\widetilde C$, $\varinjlim_\alpha\,{}_{\widetilde C}X_\alpha=a(\varinjlim_\alpha\,{}_{\widehat C}X_\alpha)$, donc $H(\varinjlim_\alpha\,{}_{\widetilde C}X_\alpha)=H(a(\varinjlim_\alpha\,{}_{\widehat C}X_\alpha)$

% L20
\[
=(H\circ a)(\mathop{\underline{\lim}}_\alpha X_\alpha)\underset{\substack{\text{par hyp}\\\text{sur }H\circ a}}{=}\mathop{\underline{\lim}}_\alpha(H\circ a)(X_\alpha)=\varinjlim_\alpha H(X_\alpha)\qquad\text{OK.}
\]

% L21–24
Donc on trouve que $\alpha'$ induit

% L21–24-right
\[
\begin{aligned}
\underline{\mathrm{Hom}}'(\widetilde C,E)\xrightarrow{\simeq}{}&
\text{sous-catégorie \inserted{pleine} de }\underline{\mathrm{Hom}}^{\unsure{F|{}}}(\widehat C,E)\text{ formée des}\\
&\text{foncteurs qui a) commutent aux }\mathop{\underline{\lim}}\text{ petites}\\
&\hphantom{\text{foncteurs qui }}\text{b) Transforme morphismes}\\
&\hphantom{\text{foncteurs qui }}\text{``bicouvrants'' de }\widehat C\text{ (i.e. tels que}\\
&\hphantom{\text{foncteurs qui }}a(u)\text{ inv.) en isom.}
\end{aligned}
\]

% L25
Donc on trouve

% L26–29
\[
\begin{aligned}
\underline{\mathrm{Hom}}'(\widetilde C,E)\xrightarrow[\simeq]{\alpha_0}{}&
\text{sous-catégorie pleine de }\underline{\mathrm{Hom}}(C,\unsure{E|\deleted{\widetilde E}E})\\
&\text{formée des }u_0\colon C\longrightarrow E\text{ tels que}\\
&\widehat u_0\colon\widehat C\longrightarrow E\text{ transforme morphismes}\\
&\text{``bicouvrants'' en isom.}
\end{aligned}
\]

% L30–32
Mais cette condition implique que \deleted{\unsure{la|tt}} famille couvrante \inserted{$X_\alpha\longrightarrow X$} de $\unsure{\overline C|C|\widetilde C}$ (définissant donc une famille épimorphique $\varepsilon_C(X_\alpha)\longrightarrow\varepsilon_C(X)$ de $\widetilde C$) \deleted{est} transformée en famille épimorphique de $E$.

% L33
Cela \uline{est-il suffisant} ? Supposons cette condition remplie,

% L33-left
\margin{\uline{non}}

% foot
% archive sheet number: 20
